% EGA V, English translation by P. Blass, 1987 (cote 24): archive cote 24, site page 4 (PDF page 5).
% typescript. Draft compared: none.
% AI-transcribed by Claude Opus 5.5 (claude-opus-5-5, adaptive thinking, effort high), Message Batches API, batch pilot (msgbatch_0189Nhf1xB615YWnztubtLfc), 2026-11-12; not yet checked by eye.
\scanpage{24/4}
\mspage{4}

\uline{HYPERPLANE SECTIONS AND CONIC PROJECTIONS}

\uline{Section 1: Preliminaries and Notation.}

Let S be a prescheme, let $\underline{E}$ be a locally free module of
finite type over S let $\underline{E}^{\vee}$ be its dual. We denote by P = $\underline{P}(\underline{E})$ the
projective fibration defined by $\underline{E}$ and by P$^{\vee}$ the projective
fibration defined by E$^{\vee}$. We shall call P$^{\vee}$ the \uline{scheme of}
\uline{hyperplanes} of P. This terminology can be justified as follows.
Let $\xi$ be a section of P$^{\vee}$ over S which is therefore determined by an
invertible quotient module L of $\underline{E}^{\vee}$. From it we obtain an
invertible quotient L$_{P}$ of $(\underline{E}^{\vee})_{\underline{P}} = (\underline{E}_{P})^{\vee}$, on the other hand, we have
the invertible quotient module O$_{P}$(1) of E$_{P}$. Passing to duals we
may take $\underline{L}_{P}{-1}$ resp. $\underline{O}_{P}(-1)$ to be invertible submodules (locally
direct \uline{summands or factors}) of $\underline{E}_{P}$ (resp. of $(\underline{E}_{P})^{\vee}$) and the pairing
$EP \otimes E^{\vee}_{P} \to \underline{O}_{P}$ defines therefore a natural pairing (*) $\underline{O}_{P}(-1) \otimes \underline{L}_{P}{-1}$
$\to \underline{O}_{P}$ or also the transpose homomorphism (**) $\underline{O}_{P} \to \underline{O}_{P}(1) \otimes \underline{L}_{P} =$
$\underline{O}_{P}(1)$ i.e. a section of $\underline{L}_{P}(1)$ canonically defined by $\xi$. The
``divisor'' of that section i.e. the closed subscheme H$_{\xi}$ of $\underline{P}$ defined
by the image ideal of (*) is called the hyperplane in $\underline{P}$ defined by
the element $\xi \in$ P$^{\vee}$(S). We could describe it by noting that locally
over $\underline{S}$, $\xi$ is given by a section $\phi$ of E such that $\phi$ (s) $\neq 0$ for all
s ($\phi$ is determined by $\xi$ up to multiplication by an invertible
section of $\underline{O}_{S}$); since E = p$^{*}$(O$_{P}$(1)) (p: P $\to$ S being the
projection), $\phi$ can be considered as a section of O$_{P}$(1) the divisor
of which is nothing else but H$_{\xi}$.

Of course, if we consider L-1 as an invertible submodule of E
locally a direct factor in E then the correspondence between $\xi$
% NOTE: a white strip covers the next typed line (only faint traces show through); nothing else is typed below it on the page
